\chapter{Complexity — Unary Tape}
%%% generated by the blueprint pipeline from TCSlib.Complexity.TuringMachine.UnaryTape
%%%%%%%%%%%%%%%%%%%%%%%%%%%%%%%%%%%%%%%%%%%%%%%%%%%%%%%%%%%%%%%%%%%%%%%%%%%%%%%
\section{Overview}

This module describes work tapes of a multi-tape Turing machine [AB09, \S 1.2] that hold a
natural number in unary. A work tape is a map from the integer cell positions to the symbols
$\{0, 1, \sqcup\}$, where $\sqcup$ is the blank. The unary tape of $m$ holds $1$ on cells
$0, \dots, m-1$ and is blank everywhere else. Machines that keep counters in unary increment a
counter by writing $1$ on its first blank cell and decrement it by erasing its last cell. The
module proves that these two writes turn the unary tape of $m$ into that of $m+1$, and the unary
tape of $j+1$ into that of $j$. It also defines the effect of an optional write at a head
position.

%%%%%%%%%%%%%%%%%%%%%%%%%%%%%%%%%%%%%%%%%%%%%%%%%%%%%%%%%%%%%%%%%%%%%%%%%%%%%%%
\section{Declarations}

\begin{definition}[Unary tape of a natural number]
\lean{Turing.UnaryTape.ones}
\label{Turing.UnaryTape.ones}
\leanok
For a natural number $m$, the unary tape of $m$ is the tape
$\mathrm{ones}_m : \mathbb{Z} \to \{0, 1, \sqcup\}$, where $\sqcup$ denotes the blank symbol,
given by
\[
  \mathrm{ones}_m(c) =
  \begin{cases}
    1 & \text{if } 0 \le c < m, \\
    \sqcup & \text{otherwise.}
  \end{cases}
\]
\end{definition}

\begin{theorem}[Unary tape holds 1 below its value]
\lean{Turing.UnaryTape.ones_lt}
\label{Turing.UnaryTape.ones_lt}
\leanok
\uses{Turing.UnaryTape.ones}
\difficulty{1}
Let $m \in \mathbb{N}$, and let $\mathrm{ones}_m : \mathbb{Z} \to \{0, 1, \sqcup\}$ be its unary
tape, which holds $1$ on cells $0, \dots, m-1$ and the blank $\sqcup$ elsewhere. If
$c \in \mathbb{Z}$ satisfies $0 \le c$ and $c < m$, then $\mathrm{ones}_m(c) = 1$.
\end{theorem}

\begin{theorem}[Unary tape is blank from its value on]
\lean{Turing.UnaryTape.ones_ge}
\label{Turing.UnaryTape.ones_ge}
\leanok
\uses{Turing.UnaryTape.ones}
\difficulty{3}
Let $m \in \mathbb{N}$, and let $\mathrm{ones}_m : \mathbb{Z} \to \{0, 1, \sqcup\}$ be its unary
tape, which holds $1$ on cells $0, \dots, m-1$ and the blank $\sqcup$ elsewhere. If
$c \in \mathbb{Z}$ satisfies $m \le c$, then $\mathrm{ones}_m(c) = \sqcup$.
\end{theorem}

\begin{theorem}[Unary tape is blank at negative cells]
\lean{Turing.UnaryTape.ones_neg}
\label{Turing.UnaryTape.ones_neg}
\leanok
\uses{Turing.UnaryTape.ones}
\difficulty{3}
Let $m \in \mathbb{N}$, and let $\mathrm{ones}_m : \mathbb{Z} \to \{0, 1, \sqcup\}$ be its unary
tape, which holds $1$ on cells $0, \dots, m-1$ and the blank $\sqcup$ elsewhere. If
$c \in \mathbb{Z}$ satisfies $c < 0$, then $\mathrm{ones}_m(c) = \sqcup$.
\end{theorem}

\begin{theorem}[Unary tape of zero is blank]
\lean{Turing.UnaryTape.ones_zero}
\label{Turing.UnaryTape.ones_zero}
\leanok
\uses{Turing.UnaryTape.ones}
\difficulty{3}
For $m \in \mathbb{N}$, let $\mathrm{ones}_m : \mathbb{Z} \to \{0, 1, \sqcup\}$ be the unary
tape of $m$, which holds $1$ on cells $0, \dots, m-1$ and the blank $\sqcup$ elsewhere. The unary
tape of $0$ is the all-blank tape: $\mathrm{ones}_0(c) = \sqcup$ for every $c \in \mathbb{Z}$.
\end{theorem}

\begin{theorem}[Incrementing a unary tape]
\lean{Turing.UnaryTape.update_ones_succ}
\label{Turing.UnaryTape.update_ones_succ}
\leanok
\uses{Turing.UnaryTape.ones}
\difficulty{3}
For $m \in \mathbb{N}$, let $\mathrm{ones}_m : \mathbb{Z} \to \{0, 1, \sqcup\}$ be the unary
tape of $m$, which holds $1$ on cells $0, \dots, m-1$ and the blank $\sqcup$ elsewhere. For every
$m \in \mathbb{N}$, the tape obtained from $\mathrm{ones}_m$ by setting cell $m$ to $1$ and leaving
every other cell unchanged is $\mathrm{ones}_{m+1}$.
\end{theorem}

\begin{theorem}[Decrementing a unary tape]
\lean{Turing.UnaryTape.update_ones_pred}
\label{Turing.UnaryTape.update_ones_pred}
\leanok
\uses{Turing.UnaryTape.ones}
\difficulty{3}
For $m \in \mathbb{N}$, let $\mathrm{ones}_m : \mathbb{Z} \to \{0, 1, \sqcup\}$ be the unary
tape of $m$, which holds $1$ on cells $0, \dots, m-1$ and the blank $\sqcup$ elsewhere. For every
$j \in \mathbb{N}$, the tape obtained from $\mathrm{ones}_{j+1}$ by setting cell $j$ to the blank
$\sqcup$ and leaving every other cell unchanged is $\mathrm{ones}_j$.
\end{theorem}

\begin{definition}[Optional write at a head position]
\lean{Turing.UnaryTape.wrT}
\label{Turing.UnaryTape.wrT}
\leanok
Let $T : \mathbb{Z} \to \{0, 1, \sqcup\}$ be a tape, where $\sqcup$ denotes the blank symbol, let
$p \in \mathbb{Z}$ be a head position, and let $w$ be an optional write, which is either
``no write'' or a symbol $s \in \{0, 1, \sqcup\}$ to be written. The tape after the write is $T$
itself if $w$ is ``no write''. If $w$ writes $s$, it is the tape that holds $s$ at cell $p$ and
agrees with $T$ at every other cell.
\end{definition}
